\section{Positive bridges, caps, and growth}\label{L:part10}

We now turn the partition estimates into positive families of open paths. The immediate goals are transverse tightness, length-controlled half-plane growth with recoverable cuts, and caps at prescribed boundary gaps. These are the geometric inputs needed to convert a long arch into a bridge without losing more than logarithmic factors. Use triangle-center walks with boundary endpoints at midports, weight \(X^{|\omega|}\) per walk (\(|\omega|\) counts centers), so gluing at a midport multiplies weights. Triangle-side levels in any fixed lattice direction are numbered at unit spacing (one interline distance as unit). A strict bridge of height \(h\ge1\) goes from a fixed port on level 0 to level \(h\), with all its centers strictly between. By Proposition~\ref{L:prop:strip-mass} its total mass is \(B_h\asymp h^{-1/4}\); include \(B_0=1\). Reflection/translation give the same conventions in any of the side directions. The arch on a half plane is a nonempty return to the source side before any other contact, and \(A_g\) denotes its mass to a port at positive horizontal gap \(g\) (one of the two displacements), integer in lattice side lengths.

We recall the useful clipping rule: for a boundary port of a finite convex domain of triangles with straight sides, exit flux with coefficient \(\cos((3/8)\,{\rm turn})\) is 1 by the extension identity. Coefficients are fixed by boundary tangents and uniformly positive. Cutting such a domain by a lattice side line (keeping the source) thus bounds the mass of all lost exits, including excursions of arcs to retained endpoints past the cut, by a constant times new exits at the cut; conversely positive lost terms give lower bounds on these new exits. Indeed coefficients to shared endpoints agree. More generally on a simple polygonal domain made of triangles with the source on its outer (counterclockwise) boundary, the same complex identity from the finite proof holds: collision detours are approached from outside so still cancel. If \(\alpha\) is the tangent angle at an exit port lifted by following this boundary from the source (angle measured from the starting boundary tangent), the path turn to that port is \(\alpha-\pi\), by closing counterclockwise from there to the source (the two extra corners are \(+\pi/2\) each). In the straight strip the identity \(1=\cos(3\pi/8) A^{[h]}+B_h\) holds as in the strip calculation, \(A^{[h]}\) summing both directions of return. Thus total half-plane arch mass is \(1/\cos(3\pi/8)\), and the arch height tail has order \(B_h\).

Here are two tightness consequences with fixed constants.
\begin{enumerate}
\item Arch mass of diameter \(>r\) is \(\le C r^{-1/4}\). Take an equilateral triangle of side order \(r\), slightly smaller than \(r\), with horizontal base. For a source in the central third of the base compare to the half-plane identity. All arches of diameter \(>r\) are lost, hence bounded by exits on the other sides. Average over base sources there and reverse from those sides: for each other-side port, arcs ending on the central base must penetrate to normal distance \(c r\) or larger from that port's side. Clipping by a parallel line at such distance thus bounds their sum by \(C B_{\lfloor c'r\rfloor}\). There are \(O(r)\) ports on the other sides, proving the bound.
\item Bridges of height \(d\) remaining at transverse displacements at most \(Td\) from the source's normal have mass \(\ge c d^{-1/4}\) uniformly for sufficiently large fixed \(T\), even specifying either weak sign for the endpoint displacement from that normal. Indeed clip the strip at lateral distance of order \(S\gg d\) by two slant lattice side lines. Mass lost is bounded by lateral exits by the strip identity. Those exits are included in the analogous half-plane clip with same lateral lines and a distant parallel top at height \(\asymp S\). There the half-plane identity bounds their sum by lost arches, all of diameter \(\gtrsim S\). This proves tightness by the preceding estimate; for a sign use reflection symmetry in the source's normal.
\end{enumerate}

The enclosing-loop estimate supplies length control for the growth
construction below. We record its finite-domain consequence before choosing
the growth prefixes, so that length can be included in their cut tests.

\begin{proposition}[Finite chord length mass]\label{L:prop:box-moment}
For every finite convex domain of triangles with a bounded number of
lattice sides and diameter at most $C_0r$, the critical length-weighted
mass of all boundary-to-boundary chords, with both endpoints summed, is
at most $C r^{25/12}$. In particular, for a fixed $T$, the length-weighted
mass of strict bridges of height $d$ from one fixed source and transverse
displacement bounded by $Td$ is at most $C_Td^{13/12}$.
\end{proposition}
\begin{proof}
By the plane ball upper bound following \eqref{L:eq1} and the visit comparison \eqref{L:eq6}, for a finite convex polygonal domain as there of diameter \(O(r)\) with boundedly many lattice sides, the sum of chordal masses times length, over all sources and exits, is \(O(r^{2+1/12})\) by the ball upper bound. Centers right at the boundary where \eqref{L:eq6} may not apply number \(O(r)\), costing only \(O(r^2)\) since total exit mass per source is bounded. In particular bridges of height \(d\) from a fixed source staying within transverse displacement \(Td\) as in tightness have length-weighted mass at most \(C_T d^{13/12}\). Indeed shift them to order \(d\) distinct sources in a common larger such box with the same two strip walls.
\end{proof}

A cut is read at a port on a lattice-side grid line. In the forward
paths considered below, each candidate \(u\)-level lies between the
starting and terminal levels. An interior candidate is a
\emph{single-crossing cut} if the path crosses it at exactly one port.
A terminal candidate is treated as a cut when all centers lie on its
starting side; it becomes single crossing after an extension strictly
beyond that wall. The \emph{first qualifying cut} is the lowest candidate
with this cut property whose prefix satisfies the specified diameter,
confinement, terminal-position, and, when present, length conditions.
These auxiliary conditions concern the prefix alone. If \(u=t\) is
selected, every center of the suffix lies in \(u>t\), or the suffix is
empty at a terminal cut. Removing that suffix or replacing it by another
extension in \(u>t\) changes neither the crossing counts at earlier
candidates nor their prefix conditions, and the selected level remains
eligible as a terminal or single-crossing cut. Thus the first-cut rule
recovers the same cut.

\begin{unnumberedproposition}[Recoverable half-plane growth]
\phantomsection\label{L:prop:recoverable-growth}
\label{L:prop:length-controlled-growth}
Let \(n,u\) be signed normal level coordinates in two lattice directions
whose unit inward vectors have dot product \(1/2\). Let \(A\) be a port
on the \(n=0\) wall, and measure both coordinates relative to \(A\).
There are constants \(r_0,c,C,c'>0\), independent of \(r\), such that
for every \(r\ge r_0\) there is a family of distinct prefixes from \(A\)
with total critical mass at least \(c\), whose terminal \(u\)-levels
\(s\) lie in the grid bin \([r,2r]\). Every prefix has all centers in
\(n>0\) and \(u<s\), diameter at most \(Cr\), center length at most
\(Cr^{4/3}\), and terminal \(n\)-coordinate at least \(c'r\).
Its terminal cut is selected by the first-cut rule and is recoverable
after any extension strictly beyond it in \(u\).

For each fixed sufficiently small \(\gamma>0\), there are a threshold
\(r_\gamma\ge r_0\) and positive constants \(K_\gamma,c_\gamma,C_\gamma\),
and \(C_{1,\gamma}\) such that, for every \(r\ge r_\gamma\) and every
admissible grid target \(S\ge K_\gamma r\), these prefixes can be
continued to the exact wall \(u=S\), with terminal port free, so that
all centers remain in \(n>0\) and \(u<S\), the diameter is at most
\(C_\gamma S\), the center length is at most \(C_\gamma S^{4/3}\),
and the resulting family has total mass at least
\[
 c_\gamma S^{-1/4}(S/r)^{-C_{1,\gamma}}.
\]
During the continuation, the transverse distance from the positive
\(u\)-normal ray through the selected prefix endpoint is at most
\(\gamma u\). The cuts separating all stages are recovered by the same
first-cut rule.
\end{unnumberedproposition}

The length restrictions will be needed for the final exterior connectors.
The cap construction below uses only the mass, geometry, and recovery
conclusions; retaining the length restrictions costs only fixed constants.

\begin{proof}
\textbf{Initial exits and the first cut.}
Fix a target grid level \(s\in[r,2r]\). Truncate the half-plane \(n>0\)
to a finite convex domain of diameter \(O(r)\) that contains the arches
of diameter at most \(C''r\) used below, and clip it at \(u=s\).
The half-plane height tail supplies arch mass at least \(cr^{-1/4}\)
reaching \(n>8r\). Choosing \(C''\) large retains a fixed fraction by
the arch diameter tightness. Reflection in the starting \(n\)-normal
ray shows that at least half of this retained mass crosses the clipping
line. The positive clipping identity therefore supplies exits at \(u=s\)
of mass at least \(cr^{-1/4}\).

We need the endpoint to be a positive distance inside \(n>0\).
Translate the domains along the \(n\)-wall so that all target \(u\)-walls
in the bin coincide. The starting ports then become distinct. On that
common wall there are \(O(1+c'r)\) ports with \(n<c'r\). Reverse paths
from each such port. They are included in a common finite convex
truncation between the two walls, and reaching the translated starting
ports costs at most \(Cr^{-1/4}\) per reversed source: clip at a line
parallel to the aligned \(u\)-wall at distance of order \(r\), then use
the bridge upper bound. Thus the total shallow-endpoint mass in the bin
is at most \(C(1+c'r)r^{-1/4}\). The exits over all levels in the bin
have mass at least \(cr^{3/4}\); choose \(c'\) small and then \(r\)
large to retain that order of good exits.

These translated exits are boundary chords in one finite convex domain
of diameter \(O(r)\), after enlarging its other sides by a fixed factor.
Proposition~\ref{L:prop:box-moment} bounds their total length-weighted
mass by \(Cr^{25/12}\). Exits of length greater than \(Mr^{4/3}\)
therefore have mass at most
\[
 (Mr^{4/3})^{-1}Cr^{25/12}=(C/M)r^{3/4}.
\]
Choose the fixed \(M\) large enough to retain a positive fraction of
the good exits. Undo the translations, so that they again start at \(A\).

For each retained exit, keep the prefix to its first qualifying
\(u\)-cut in the bin. The tests are the diameter bound, \(n>0\)
confinement, terminal-depth bound, and length at most \(Mr^{4/3}\).
All concern the prefix alone, and the final wall qualifies. After a
candidate single crossing, the discarded suffix is a strict
\(u\)-bridge, and its total possible mass over the target levels is
at most
\[
 \sum_{0\le d\le2r}B_d=O(r^{3/4}).
\]
Consequently the distinct selected prefixes have total mass at least
a positive constant. Their lengths are at most \(Mr^{4/3}\).
The selected prefix and every earlier test are unchanged on extension
strictly beyond its cut, so that cut remains recoverable.

\smallskip\noindent\textbf{Length-controlled increments.}
Fix the transverse confinement constant \(T\) from the tightness
estimate. A strict bridge of increment \(d\) within that confinement
has mass at least \(cd^{-1/4}\), even with either specified weak
terminal sign. The whole confined family has length-weighted mass
at most \(C_Td^{13/12}\), by Proposition~\ref{L:prop:box-moment}.
Thus imposing length at most \(M_1d^{4/3}\) loses at most
\((C_T/M_1)d^{-1/4}\) from either sign-restricted family. Choose
\(M_1\) once, large enough for the confinement and sign constants.
Every increment family below then retains its original mass order.

Now start from a selected cut at \(u=t\). Choose a small fixed
\(\zeta>0\), with its value in terms of \(\gamma\) fixed below.
After fixing \(\zeta\), take \(r_\gamma\) large enough that the integral
increment bins below have their asserted order of cardinality and any
fixed lattice-scale rounding errors are absorbed by the transverse bound.
While \(t\le S/(1+3\zeta)\), use strict \(u\)-bridges
of integral increment \(d\in[\zeta t,2\zeta t]\). Require transverse
displacement at most \(Td\) from the increment's starting normal and,
if the current endpoint is off the initial \(u\)-normal ray, require
the weak sign pointing toward that ray. With the length restriction
just imposed, their mass over the bin is at least
\(c_\zeta t^{3/4}\).

Extract first qualifying cuts in this increment bin, allowing the
terminal wall and testing
transverse displacement at most \(2T\zeta t\), the selected terminal
sign, and prefix length at most \(M_1(2\zeta t)^{4/3}\).
These tests depend on the current starting cut and its prefix, not on
the discarded terminal target \(d\). Every retained increment has
its terminal wall eligible. The suffix sum is again
\(O_\zeta(t^{3/4})\), so the distinct increment prefixes have mass
at least \(c'_\zeta>0\), and their cuts are recoverable.

Each such stage increases \(t\) by at least \(\zeta t\), so there are
\(O_\zeta(1+\log(S/r))\) stages. At the first larger \(t\), the
remaining increment \(S-t\) is comparable to \(\zeta S\); finish
with a strict bridge at that exact height, with the same confinement,
weak terminal sign, and length restriction. It has mass at least
\(c_\zeta S^{-1/4}\). The weak signs keep endpoint errors bounded
by the largest increment-width bound used so far; within a stage the
additional error has the same order. Thus the continuation stays
within \(O(T\zeta u)\) of the initial ray. That ray begins at
\(n\ge c'r\) and gains \(n\)-clearance at slope \(1/2\).
Choose \(\zeta\) small enough in terms of \(\gamma,c',T\).
All stages then remain in \(n>0\), the diameter is \(O(S)\),
and strict slabs keep new centers beyond the preceding cuts and
before the target wall.

The starting levels of successive stages grow geometrically.
Their prefix-length bounds therefore sum to \(O_\gamma(S^{4/3})\).
The final increment is comparable to \(\zeta S\), so its length and
the initial \(O(r^{4/3})\) contribution obey the same bound.
Critical weights multiply at the joining ports. The first-cut rule
recovers the stages from the output, so multiplying their lower bounds
does not create an unaccounted multiplicity. A fixed positive factor
for each of \(O_\zeta(1+\log(S/r))\) stages, followed by the exact
final bridge, is at least
\(c_\gamma S^{-1/4}(S/r)^{-C_{1,\gamma}}\).
This proves every asserted conclusion.
\end{proof}

We next construct caps at prescribed gaps.  The first step gives a
lower bound at some gap on each scale.  A one-cell-thick boundary
deformation then compares any smaller prescribed gap to that one.

\begin{unnumberedproposition}[Caps at prescribed gaps]
\phantomsection\label{L:prop:caps}
There are constants $c,C>0$ such that, for every integer $g\ge1$,
\[
 \sum_{\substack{\omega\text{ an arch at gap }g\\
                  \operatorname{diam}\omega\le Cg}}
       \mu(\omega)\ge c g^{-5/4}.
\]
Moreover, the unrestricted arch masses satisfy
$A_g\ge c A_q$ whenever $1\le g<q$.
\end{unnumberedproposition}
\begin{proof}
\medskip\noindent\textbf{A gap with substantial mass on each scale.}
We first show that, for every sufficiently large $r$, some integer
$q'\in[\lfloor r\rfloor,Cr]$ has arch mass at least $cr^{-5/4}$
with diameter at most $Cr$.
Use level coordinates $Y,U_+,U_-$ whose normals have angles
$\pi/2,\pi/6,-\pi/6$, respectively.  With a lattice vertex as origin
and one interline distance as unit, $U_-=U_+-Y$, and the integer levels
are lattice-side grid lines.  Put the source at the midpoint of
$[0,1]$ on $Y=0$, using horizontal lattice-side length one.  Set
$p=\lfloor r\rfloor$.  Choose integers $S,M$ that are successively
large fixed multiples of $r$, up to rounding, and let $D_r$ be the
finite triangle domain
\[
 Y\le M,\qquad -M\le U_-\le S,\qquad
 (Y\ge0\ \text{or}\ U_+\ge p).
\]
Its lower boundary travels east to position $p$, then down-right
along $U_+=p$ to the wall $U_-=S$.  That wall rises up-right to the
top.  Clipping by $Y\ge0$ replaces this lower excursion by the
horizontal segment from $p$ to $S$.

Let $\alpha$ be the tangent angle lifted along the counterclockwise
outer boundary from the source, as in the boundary-flux identity.
In $D_r$ its values lie in $[-\pi/3,2\pi]$; the clipped domain has
the same lifts at shared boundary ports.  Rotate the complex identity
by $\exp(i(3/8)\pi/6)$ and take real parts.  The exit coefficient is
\[
 \cos\!\left(\frac38\left(\alpha-\frac{5\pi}{6}\right)\right)
 \ge\cos(7\pi/16)>0
\]
throughout both domains.  Consequently the mass of exits removed by
the clipping gives a lower bound, up to a fixed factor, for the mass
of exits on its new horizontal segment.  Indeed shared coefficients
agree, and enlarging the domain only increases the path mass at a
shared exit.

We produce removed exits of mass at least $cr^{-1/4}$.  Apply the
initial-prefix part of recoverable half-plane growth with normals
$(Y,U_+)$ and scale $2r$.  The resulting prefixes have constant total
mass, diameter $O(r)$, and terminal absolute $U_+$-level greater
than $p$.  From each terminal port apply the full growth construction
with normals $(U_+,U_-)$, initial scale $r$, and exact absolute target
$U_-=S$.  All new centers lie strictly beyond the first cut in $U_+$.
They are therefore disjoint from the first prefix and remain in the
additional part of $D_r$ even when $Y<0$.  The ratio $S/r$ is fixed,
so the second construction has mass at least $cr^{-1/4}$.

Along a positive $U_-$-normal ray, $Y$ decreases at rate $1/2$.
The initial stages contribute an $O(r)$ displacement and the
continuation contributes transverse error $O(\gamma S)$.
Choose $\gamma$ small, then $S/r$ large enough that the terminal
port lies below $Y=0$.  Finally choose $M/r$ still larger so that
all paths fit below the top and to the right of the other wall.
The exact-target property keeps all centers before $U_-=S$.
The first $U_+$-cut is recovered from the output because the second
part lies strictly beyond it.  Hence multiplying the two family
weights introduces no additional multiplicity.  We have obtained
the claimed mass of exits on the removed part of $U_-=S$.

Subtracting the two rotated flux identities now gives at least
$cr^{-1/4}$ mass on the new horizontal segment of the clipped domain.
These paths are arches, have diameter $O(r)$, and end at only $O(r)$
possible gaps in $[p,S]$.  One gap $q'$ therefore has the asserted
mass $cr^{-5/4}$.

\medskip\noindent\textbf{Gap comparison.}
\phantomsection\label{L:cap-gap-comparison}
Fix integers $1\le g<q$, put $p=g+1$, and let $D$ be a large finite
convex triangle domain above the horizontal base.  Its base contains
the source $1/2$ and the target $q+1/2$.  Write $A_q(D)$ for the
mass of arches between these ports contained in $D$.  In complex
coordinates put $w=e^{-i\pi/3}$ and add the one-cell-thick bump
\[
 [p,q+1]+[0,w]
\]
below the base.  Equivalently, add the rhombi
$j+[0,1]+[0,w]$, $j=p,\ldots,q$.  Call the enlarged domain $D'$.

Figure~\ref{L:fig:cap-comparison} shows the added cells and the fixed
extensions used below.
\begin{figure}[htbp]
\centering\begin{tikzpicture}[x=.8cm,y=.8cm,>=Stealth,font=\small]
 \fill[blue!4] (-.4,0)--(9.8,0)--(9.8,1.8)--(-.4,1.8)--cycle;
 \node[blue!65!black] at (4.8,1.25) {$D$};
 \draw[gray] (-.4,0)--(9.8,0);
 \fill[red!9] (4,0)--(9,0)--(9.5,-.866)--(4.5,-.866)--cycle;
 \draw[thick] (4,0)--(4.5,-.866)--(9.5,-.866)--(9,0);
 \foreach \j in {5,6,7,8} {
   \draw[gray!60] (\j,0)--({\j+.5},-.866);
 }
 \foreach \j in {0,3} {
   \draw[dashed,blue!60!black] (\j,0)--({\j+1},0)--({\j+1.5},-.866)--({\j+.5},-.866)--cycle;
 }
 \draw[very thick,blue!65!black] (.5,0)--(.5,-.289)--(1,-.577)--(1,-.866);
 \draw[very thick,blue!65!black] (4.25,-.433)--(4,-.577)--(4,-.866);
 \draw[very thick,red!65!black] (8.5,0)--(8.5,-.289)--(9,-.577)--(9.25,-.433);
 \fill (.5,0) circle(1.8pt) node[above] {$1/2$};
 \fill (8.5,0) circle(1.8pt) node[above left] {$q+1/2$};
 \fill (4.25,-.433) circle(1.8pt) node[above right] {$F_{\rm L}$};
 \fill (9.25,-.433) circle(1.8pt) node[right] {$F_{\rm R}$};
 \fill (1,-.866) circle(1.8pt);
 \fill (4,-.866) circle(1.8pt);
 \draw[dotted] (1,-.94)--(1,-1.55);
 \draw[dotted] (4,-.94)--(4,-1.55);
 \draw[<->] (1,-1.4)--node[below] {$g=p-1$}(4,-1.4);
 \node[fill=red!9,inner sep=2pt] at (6.7,-.43) {added bump};
 \node[above] at (4,0) {$p$};
 \node[above right] at (9,0) {$q+1$};
\end{tikzpicture}

\caption{Local geometry of the gap comparison, shown for $p=4$ and $q=8$.
The shaded bump enlarges $D$ to $D'$. The red extension sends a path
ending at $q+1/2$ to the right slant port. The two blue extensions use
the dashed cells outside $D'$ to turn a path to the left slant port
into an arch of gap $g=p-1$. The labels $F_{\rm L},F_{\rm R}$ denote
exit masses at the slant ports.}\label{L:fig:cap-comparison}
\end{figure}

Use the complex flux identity at the same source in $D$ and $D'$,
multiply it by $e^{i(3/8)\pi}$, and take imaginary parts.  An exit
with lifted tangent angle $\alpha$ now has coefficient
$\sin(3\alpha/8)$.  On the left slant of the bump,
$\alpha=-\pi/3$, giving $-\sin(\pi/8)$; on its right slant,
$\alpha=2\pi/3$, giving $\sin(\pi/4)$.  The bottom horizontal
ports and the remaining horizontal ports to the right of the source
have coefficient zero.  At every other shared boundary port the
coefficient is unchanged and nonnegative, since the convex boundary
lifts lie in $[0,2\pi]$.

Let $F_{\rm L}$ and $F_{\rm R}$ be the masses of paths from the
source to the left and right slant ports of the bump in $D'$.
Subtract the identity in $D$ from the one in $D'$.  The right-hand
sides agree, and the mass at a shared port can only increase.  Thus
\[
 \sin(\pi/8)F_{\rm L}\ge\sin(\pi/4)F_{\rm R}.
\]
Every arch counted by $A_q(D)$ can be extended through the last bump
cell, cell $q$, to the right slant port.  This adds a bounded number
of centers in a previously unused region and is recoverable by
deleting that fixed extension.  Hence $F_{\rm R}\ge cA_q(D)$.

Conversely, extend a path counted by $F_{\rm L}$ through cell $p-1$
to that cell's bottom port, and extend its source backwards through
cell $0$ to the bottom port there.  Both cells lie outside $D'$ and
are distinct because $p\ge2$, so these extensions avoid the path
and each other.  The new endpoints are
\[
 \frac12+w\qquad\text{and}\qquad p-\frac12+w.
\]
Their gap is $p-1=g$.  All centers are strictly above the common
bottom line; translating by $-w$ gives an arch at gap $g$ in the
original half-plane.  The fixed extensions have bounded weight cost
and bounded inverse multiplicity.  Therefore
\[
 A_g\ge cF_{\rm L}\ge cA_q(D).
\]
The constants are independent of $D,g,q$.  Exhausting the half-plane
by such domains proves the unrestricted comparison $A_g\ge cA_q$.
This comparison applies to all paths at the farther gap, without a
diameter restriction.  If $q=O(r)$ and the input arches are contained
in a domain of diameter $O(r)$, the bump and both extensions also
have diameter $O(r)$, so the comparison preserves such a cutoff up
to a fixed factor.

For large $g$, apply the first part with $r=g+1$.  It gives
$q'\ge g+1$ with mass at least $cg^{-5/4}$ and diameter $O(g)$.
Use the finite-domain version of the gap comparison to obtain the
same lower bound at $g$, still with diameter $O(g)$.  The finitely
many smaller gaps follow by choosing one sufficiently large fixed
$r$ and adjusting $c,C$.  This proves the proposition.
\end{proof}

The same comparison gives the pointwise upper bound. For large $g$,
each of order $g$ integer gaps $j\in[g/2,g)$ satisfies $A_j\ge cA_g$.
Every arch at such a gap has diameter at least $g/2$, so the diameter
tail proved above gives
\[
 gA_g\le C\sum_{g/2\le j<g}A_j\le Cg^{-1/4}.
\]
Thus $A_g\le C(1+g)^{-5/4}$, with bounded gaps absorbed by changing $C$.

\subsection{Critical irreducible bridges}\label{L:step17}

As a separate consequence of the finite bridge bounds, we obtain the
classical irreducible-bridge law. The following result is not needed for
the logarithmic window proved in the next section. We use the classical
construction \cite{Kesten1963,MadrasSlade1993}. Its critical normalization
and independent concatenation in the honeycomb mid-edge convention are
also given in \cite[Appendix]{BeatonEtAl2014}; the finite bridge bounds
here yield the additional quantitative tails below.
Cutting a strict bridge at all its intermediate single-crossing levels
decomposes it uniquely into strict bridges with no intermediate
single-crossing level; these are the \emph{irreducible bridges}. The cuts are ordered by height.
Each portion lies strictly between its two cut lines: a return past a
line would cross it again. Conversely, irreducible pieces concatenate
freely, translating the next starting port to the previous terminal
port. Their center sets lie in disjoint open slabs, so the concatenation
is self-avoiding and its critical weight is the product of the weights.

\begin{proposition}[Critical irreducible law]\label{L:prop:irreducible}
The sum of $X^{|\gamma|}$ over irreducible strict bridges $\gamma$ from a
fixed port is one. Under this probability law, if $Y$ and $D$ are the
height and diameter of the sampled bridge, then
\[
 1-\mathbb E e^{-tY}\asymp t^{3/4}\quad(0<t\le1),\qquad
 \Pr(D>r)\le C r^{-3/4}\quad(r\ge1).
\]
Independent samples may be concatenated to form an infinite
self-avoiding port path.
\end{proposition}
\begin{proof}
Let $i_j$ be the total irreducible mass at height $j\ge1$. Unique
factorization gives the identity of nonnegative power series
\[
 \sum_{h\ge0}B_hz^h=\sum_{k\ge0}\left(\sum_{j\ge1}i_jz^j\right)^k.
\]
For $z=e^{-t}<1$ the left side is finite, since $B_h\asymp h^{-1/4}$.
Consequently $I(t)=\sum i_j e^{-tj}<1$ and
\[
 1-I(t)=\left(\sum_{h\ge0}B_he^{-th}\right)^{-1}\asymp t^{3/4}.
\]
The last comparison follows by comparison of the sum with the integral
of $x^{-1/4}e^{-tx}$, separating the bounded initial terms. Monotone
convergence as $t\downarrow0$ shows $\sum i_j=1$ and proves the height
claim.

For the diameter claim, fix a sufficiently large transverse confinement
constant $T$ from the tightness argument above. Choose $c_0>0$ so small
that a bridge of height between $c_0r$ and $2c_0r$ with that confinement
has diameter at most $r$. Such bridges, summed over their heights, have
mass at least $c r^{3/4}$ for large $r$. Every irreducible piece of each
of them also has diameter at most $r$. Put $a_r=\Pr(D\le r)$. Dropping
all height restrictions and summing all concatenations of pieces with
diameter at most $r$ therefore gives
\[
 c r^{3/4}\le\sum_{k\ge0}a_r^k.
\]
If $a_r=1$ the desired tail bound is immediate. Otherwise the right side
is $(1-a_r)^{-1}$, which proves the bound. Enlarging its constant handles
bounded $r$. Finally all piece heights are positive integers, so an
infinite independent concatenation escapes through successive disjoint
slabs and defines a self-avoiding path.
\end{proof}

